\documentclass[11pt]{article}

\usepackage[T1]{fontenc}
\usepackage{lmodern}
\usepackage[a4paper,margin=27mm]{geometry}
\usepackage{amsmath,amssymb,amsthm,mathtools}
\usepackage{microtype}
\usepackage{enumitem}
\usepackage{booktabs,array,longtable}
\usepackage{xcolor}
\usepackage[colorlinks=true,linkcolor=blue!48!black,
  citecolor=blue!48!black,urlcolor=blue!55!black]{hyperref}
\setlength{\emergencystretch}{2em}

\newtheorem{theorem}{Theorem}[section]
\newtheorem{proposition}[theorem]{Proposition}
\newtheorem{corollary}[theorem]{Corollary}
\newtheorem{lemma}[theorem]{Lemma}
\theoremstyle{definition}
\newtheorem{definition}[theorem]{Definition}
\theoremstyle{remark}
\newtheorem{remark}[theorem]{Remark}
\numberwithin{equation}{section}

\newcommand{\C}{\mathbb C}
\newcommand{\Q}{\mathbb Q}
\newcommand{\R}{\mathbb R}
\newcommand{\N}{\mathbb N}
\newcommand{\cC}{\mathcal C}
\newcommand{\cG}{\mathcal G}
\newcommand{\cH}{\mathcal H}
\newcommand{\cK}{\mathcal K}
\newcommand{\cP}{\mathcal P}
\newcommand{\cR}{\mathcal R}
\newcommand{\cS}{\mathcal S}
\newcommand{\Hom}{\operatorname{Hom}}
\newcommand{\im}{\operatorname{im}}
\newcommand{\rad}{\operatorname{rad}}
\newcommand{\Trans}{\operatorname{Trans}}
\newcommand{\Rea}{\operatorname{Re}}

\title{\textbf{Relation-kernel transgression and the Riemann boundary}\\
\large A pre-spectral arithmetic route to horizontal purity}
\author{Michael Haufschild}
\date{}

\begin{document}

\maketitle

\begin{abstract}
We isolate a general conservation law for cohomological information made
exact by a cover.  Let $\pi:R\twoheadrightarrow A$ be a surjective group
map, let $K=\ker\pi$, and let $c$ be a normalized two-cocycle on $A$ with
values in a rational vector space.  If $\pi^*c=\delta b$, then $b|_K$ is a
character.  If every degree-one character of $R$ descends to $A$, this
relation-kernel character is independent of $b$ and vanishes exactly when
$c$ is already a coboundary on $A$.  For an abelian $A$ and alternating
$c$, it vanishes exactly when $c=0$.  Thus an exactifying cover does not
erase degree-two information: it lowers that information by one degree and
stores it on the relations killed by the cover.

We apply this law to a concrete arithmetic source.  At each prime, the
Calder\'on defect line and rational-Witt packet space carry an alternating
vector-valued cup with zero radical.  On a free rational cover the cup is
exact, while its restriction to the relation kernel is a canonical nonzero
character.  This gives a divisor-independent, unsquared arithmetic
calibrator.

We then organize the current Riemann route around that calibrator.  A literal
two-line prime--Gamma Cauchy source, its harmonic Hodge primitive, completed
Green transport, adaptive M\"obius--theta approximation, and resolved
provenance are already available without a critical-line premise.  Their
remaining horizontal defect is coordinatewise
\[
             g_G(\rho)=D_\rho u_\rho,
        \qquad D_\rho=2(\Rea\rho-\tfrac12),\qquad u_\rho\ne0.
\]
The codiagonal kills $(g_G,-g_G)$ identically, so no scalar trace,
positivity argument, or spectral realization can recover its sign.  We
state the one missing comparison theorem precisely: a source-defined chain
map must identify the transgression of completed Green provenance with
$D_\rho$ times the nonzero relation-kernel calibrator, while carrying the
literal Meyer--Poisson boundary to zero.  That theorem would force
$D_\rho=0$ for every zero and hence prove the Riemann hypothesis.  It is not
proved here.  The contribution is the universal transgression law, its
nonzero arithmetic realization, a closed audit of the source and transport
chain, and an exact separation of this route from Hilbert--P\'olya,
positivity, trace, and quotient-spectral programs.
\end{abstract}

\section{The decisive change of problem}

Write
\[
 \xi(s)=\frac12s(s-1)\pi^{-s/2}\Gamma(s/2)\zeta(s),
 \qquad D_\rho=2(\Rea\rho-\tfrac12)
\]
for a nontrivial zero $\rho$ of $\xi$, counted with multiplicity.  The
Riemann hypothesis is the atomic assertion $D_\rho=0$ for every such
$\rho$ \cite{Riemann1859,Titchmarsh1986}.  An identity involving
$D_\rho^2$, a sum of several $D_\rho$, or a
reflection-invariant quotient is weaker as a proof mechanism: it has already
discarded either orientation, provenance, or both.

The source problem has changed.  Earlier formulations could reasonably ask
for a source, a completed carrier, an oriented primitive, or a landing on
the Xi divisor.  Those objects now exist in a mutually compatible chain:
\begin{equation}
 \boxed{
 \begin{matrix}
 \text{prime--Gamma Cauchy boundary}
   &\longrightarrow& h_{\rm ref}
   &\xrightarrow{\ i\operatorname{sign}(D)\ }& u_G\\
 && \rotatebox{90}{$\in$} && \rotatebox{90}{$\in$}\\[-2mm]
 && \cH_\Xi &&& \cH_\Xi
 \end{matrix}}
 \label{eq:source-chain}
\end{equation}
and an independent adaptive M\"obius--theta family converges to the same
completed primitive and Green detector.  Incidence maps retain the
multiplicity copies, and a provenance-enriched quotient retains the
antidiagonal defect that direct codiagonalization erases.

The unresolved issue is therefore not representation.  It is an arithmetic
reason why the retained oriented coefficient must be zero.  This paper
proposes that the missing reason lives one degree before every spectral or
positive consumer: on the relation kernel of an exactifying arithmetic
cover.

\subsection{Result status}

Three logical levels will be kept separate.
\begin{enumerate}[label=(\roman*)]
\item The relation-kernel transgression theorem and its alternating
Calder\'on--Witt specialization are unconditional.
\item Source construction, harmonic primitive, completed Green transport,
resolved provenance, and the factorization $g_G(\rho)=D_\rho u_\rho$ with
$u_\rho\ne0$ are unconditional.
\item The chain comparison between these two unconditional structures is
open.  A conditional capstone records exactly why that one comparison would
imply RH.
\end{enumerate}
No spectral operator, positivity premise, zero enumeration, or RH-equivalent
boundedness assumption is used to obtain the arithmetic calibrator.

\section{The relation-kernel memory law}

Let $R$ and $A$ be groups, $M$ an additive abelian group, and
$\pi:R\twoheadrightarrow A$ a surjective homomorphism.  A normalized
two-cocycle is a map $c:A\times A\to M$ satisfying
\begin{align}
 c(a,d)+c(ad,e)&=c(d,e)+c(a,de),\label{eq:cocycle}\\
 c(1,a)=c(a,1)&=0.\label{eq:normalized}
\end{align}
Our sign convention declares $b:R\to M$ a primitive of $\pi^*c$ when
\begin{equation}
 c(\pi g,\pi h)=b(gh)-b(g)-b(h).
 \label{eq:primitive}
\end{equation}

\begin{theorem}[Relation-kernel transgression]
\label{thm:kernel-transgression}
Let $K=\ker\pi$ and suppose \eqref{eq:primitive} holds.
\begin{enumerate}[label=(\alph*)]
\item The restriction
\[
                  \tau_{c,\pi,b}:K\longrightarrow M,
                  \qquad \tau_{c,\pi,b}(k)=b(k),
\]
is a group homomorphism.
\item Suppose every character $\psi:R\to M$ vanishes on $K$; equivalently,
every degree-one character of $R$ descends through $\pi$.  Then
$\tau_{c,\pi,b}$ is independent of the primitive $b$.
\item Under the same degree-one hypothesis,
\[
        \tau_{c,\pi,b}=0
        \quad\Longleftrightarrow\quad
        c\text{ is a coboundary on }A.
\]
\end{enumerate}
\end{theorem}

\begin{proof}
Putting $g=h=1$ in \eqref{eq:primitive} gives $b(1)=0$.  If $k,l\in K$,
normalization gives
\[
 0=c(1,1)=b(kl)-b(k)-b(l),
\]
which proves (a).

If $b$ and $b'$ are two primitives, subtraction of their primitive
identities shows that $b-b'$ is a character on $R$.  The degree-one
hypothesis makes this character zero on $K$, proving (b).

If $c=\delta d$ downstairs, then $b=d\circ\pi$ is a primitive and vanishes
on $K$; (b) makes every transgression vanish.  Conversely suppose
$b|_K=0$.  Choose one lift $r_a\in R$ for each $a\in A$ and put
$d(a)=b(r_a)$.  Any two lifts differ by an element of $K$.  Applying
\eqref{eq:primitive} to a lift and that kernel element shows that $b$ is
constant on each fibre of $\pi$, so $d$ is well defined.  Applying
\eqref{eq:primitive} to $r_a,r_d$ and comparing $r_ar_d$ with a lift of
$ad$ gives
\[
                  c(a,d)=d(ad)-d(a)-d(d).
\]
Thus $c$ is a coboundary on $A$.
\end{proof}

The theorem is a conservation law rather than a vanishing theorem.  The
cover may make $c$ exact, but exactness records the class as a degree-one
character on precisely the relations killed by the cover.  The hypothesis
on degree one prevents a change of primitive from changing that record.

\begin{corollary}[Alternating rigidity]
\label{cor:alternating-rigidity}
Under the hypotheses of Theorem~\ref{thm:kernel-transgression}, assume in
addition that $A$ is abelian, $M$ is a vector space over $\Q$, and
$c(a,d)=-c(d,a)$.  Then
\[
               \tau_{c,\pi}=0
        \quad\Longleftrightarrow\quad c=0.
\]
\end{corollary}

\begin{proof}
A coboundary on an abelian group is symmetric.  If it is also alternating,
then $2c(a,d)=0$ for every $a,d$.  Rational divisibility gives $c=0$.
The converse is immediate from Theorem~\ref{thm:kernel-transgression}.
\end{proof}

\begin{remark}[Degree lowering, not information loss]
The law has the same logical shape as several apparently unrelated
obstructions.  A quotient remembers lost degree-zero information in its
kernel; a resolved antidiagonal remembers information killed by
augmentation; a boundary completion remembers finite-range failure in a
closure quotient.  Theorem~\ref{thm:kernel-transgression} is the
cohomological version: degree-two information made exact upstairs survives
in degree one on the relation kernel.  Only this cohomological version is
used in the proposed RH route.
\end{remark}

\section{A nonzero arithmetic calibrator}

Fix a prime $p$.  Let $Q_p$ denote the Calder\'on defect line, let $P_p$ be
the rational-Witt packet space on its affine gauge limit, and let $S_p$ be
the first-suspension packet space.  The actual source cross-effect is a
bilinear map
\[
                         B_p:Q_p\times P_p\longrightarrow S_p.
\]
More explicitly, if $\kappa_p:Q_p\to\C$ is the Calder\'on quotient
coordinate, $\Sigma_p:P_p\to S_p$ is the injective first-suspension mode,
and $h_p=p^{-1/2}$ is the polarized half-density, then
\begin{equation}
                         B_p(q,a)=\kappa_p(q)h_p\Sigma_p(a).
 \label{eq:cross-effect}
\end{equation}
Set $V_p=Q_p\oplus P_p$ and define
\begin{equation}
 \omega_p\bigl((q,a),(q',a')\bigr)
       =B_p(q,a')-B_p(q',a).
 \label{eq:CW-cup}
\end{equation}
This is not a scalar pairing.  It retains the first-suspension packet and
therefore keeps information that a trace can erase.

\begin{theorem}[Calder\'on--Witt source cup]
\label{thm:CW-cup}
For every prime $p$, $\omega_p$ is complex bilinear, alternating, and has
zero left radical:
\[
 \omega_p(x,y)=-\omega_p(y,x),\qquad
 \bigl(\forall y,\ \omega_p(x,y)=0\bigr)\Longrightarrow x=0.
\]
It is preserved literally by every gauge-corrected finite rational-Witt
refinement.  The algebraic all-prime direct sum retains each local cup in
its prime component.
\end{theorem}

\begin{proof}
Alternation follows directly from \eqref{eq:CW-cup}.  For nondegeneracy,
test the $Q_p$ component of $x$ against a fixed nonzero constant packet in
$P_p$; injectivity of the Calder\'on cross-effect forces that component to
vanish.  Test the remaining $P_p$ component against the unit Calder\'on
defect; injectivity of the first suspension forces the packet component to
vanish.  Gauge-corrected refinement commutes with both source legs of
$B_p$, so it preserves their alternating difference.  Direct-sum component
projection recovers the local identity prime by prime.
\end{proof}

View the additive group of $V_p$ multiplicatively and regard $\omega_p$ as
a normalized two-cocycle.  Choose a rational basis of $V_p$ and let
\[
                       \pi_p:R_p\twoheadrightarrow V_p
\]
be its free rational cover.  This cover has the two exact properties needed
by Theorem~\ref{thm:kernel-transgression}: rational normalized
two-cocycles pull back to coboundaries, and every degree-one character of
$R_p$ descends to $V_p$.

\begin{proposition}[Free-rational exactification]
\label{prop:free-rational}
Let $V$ and $M$ be rational vector spaces, choose a rational basis of $V$,
and let $R\twoheadrightarrow V$ be the rationalization of the free group on
that basis.  Then every normalized $M$-valued two-cocycle on $R$ is a
coboundary.  Moreover, every character $R\to M$ factors uniquely through
$V$.
\end{proposition}

\begin{proof}
For a normalized cocycle $c$ on $R$, form the central extension
$M\times_cR$ with product
\[
                 (m,g)(n,h)=(m+n+c(g,h),gh).
\]
Its kernel $M$ is uniquely divisible.  A map from the generating free group
lifts to the extension, and the universal property of rationalization
extends that lift across $R$.  The resulting section of
$M\times_cR\to R$ is a group map; its $M$ coordinate is a primitive of
$c$.  This proves degree-two exactness.  For degree one, a character on the
free rationalization is determined by the rational values of the basis
generators, hence factors through the rational vector group they span.
Surjectivity gives uniqueness.
\end{proof}

\begin{theorem}[Nonzero relation-kernel calibrator]
\label{thm:calibrator}
Let $K_p=\ker\pi_p$.  There exists a primitive $b_p:R_p\to S_p$ with
\[
       \pi_p^*\omega_p=\delta b_p
\]
such that
\begin{equation}
                \tau_p:=b_p|_{K_p}\in\Hom(K_p,S_p)
                \quad\text{is canonical and nonzero.}
 \label{eq:tau}
\end{equation}
\end{theorem}

\begin{proof}
Degree-two exactness of the free rational cover supplies $b_p$.
Degree-one descent and Theorem~\ref{thm:kernel-transgression} make its
kernel restriction independent of the chosen primitive.  If $\tau_p=0$,
Corollary~\ref{cor:alternating-rigidity} would give $\omega_p=0$, contrary
to the zero-radical statement of Theorem~\ref{thm:CW-cup} and the explicit
unit packet.  Hence $\tau_p\ne0$.
\end{proof}

This theorem supplies what positivity and spectral realization do not: a
nonzero arithmetic reference orientation defined without choosing a zero
of $\xi$.  Its nonvanishing is source-theoretic, not inferred from a
hypothetical off-line divisor.

\section{The completed Riemann source is no longer missing}

We summarize the established source chain in the notation needed for the
comparison theorem.  These results are logically independent of
Theorem~\ref{thm:calibrator}, but their meeting point is now precise.

\subsection{Literal boundary landing}

Let $z_q$ be the centered coordinate of a multiplicity copy $q$ of a Xi
zero, and fix an even reflected Gaussian
\[
                         H(z)=e^{-t_0z^2},\qquad t_0>0.
\]
On the two exterior horizontal lines $z=x\pm i$, define Cauchy columns
with the positive rectangle orientation.  Their normalized difference is
an $\ell^2$-valued boundary field $C_H(x)$.  It is Bochner integrable and
obeys the exact landing identity
\begin{equation}
                         \int_\R C_H(x)\,dx=h_{\rm ref},
              \qquad h_{\rm ref}(q)=H(z_q).
 \label{eq:cauchy-landing}
\end{equation}
Every coordinate of $h_{\rm ref}$ is nonzero.  The construction uses fixed
exterior lines and a fixed Gaussian; it contains no divisor-dependent
damping or interpolation.

\subsection{Harmonic primitive and completed Green transport}

On the multiplicity-exact Xi carrier $\cH_\Xi$, let
\[
                         W=i\operatorname{sign}(D)
\]
be the physical harmonic multiplier.  It is an isometry and $W^2=-I$.
The unique completed Green primitive of $h_{\rm ref}$ is
\begin{equation}
                             u_G=iWh_{\rm ref}.
 \label{eq:green-primitive}
\end{equation}
If harmonic orientation is applied to the literal Cauchy boundary before
integration, its integral is exactly $u_G$.  Completed Green transport
removes that orientation and reverses the real Cauchy parameter.  Hence
boundary landing, Hodge orientation, and Green transport form a commuting
triangle, rather than three representations joined only by their zero sets.

The completion is forced.  Compact source Green range is exactly the
finite-jet subspace; $h_{\rm ref}$ lies in its closure but not its range.
The Hilbert middle extends finite Green reflection uniquely to a surjective
isometry.  Thus regular finite-range realization fails, but the nonlocal
completion and its primitive are already determined.

\subsection{Independent adaptive source}

For every retained prime label, cofinal adaptive M\"obius--theta columns
converge in $\cH_\Xi$ to the same Gaussian horizon.  Harmonic orientation
converts them into primitive columns converging to $u_G$; their Green
fractures converge to the same global detector obtained from the Cauchy
boundary.  This gives two independent source presentations:
\[
 \begin{matrix}
 \text{adaptive M\"obius--theta columns}&\longrightarrow&u_G\\
 &&\rotatebox{90}{$=$}\\[-1mm]
 \text{oriented prime--Gamma Cauchy boundary}&\xrightarrow{\int}&u_G.
 \end{matrix}
 \label{eq:common-boundary}
\]
Their equality closes source existence and Green transport.  It does not
prove that the resulting antidiagonal belongs to a conservative arithmetic
sector; that is exactly where the calibrator comparison must act.

\section{The retained horizontal obstruction}

Completed Green covariance produces one vector $g_G\in\cH_\Xi$.  At a
divisor copy $q$ over $\rho_q$ it has the factorization
\begin{equation}
             g_G(q)=D_{\rho_q}\,u_q,
       \qquad D_{\rho_q}=2(\Rea\rho_q-\tfrac12),
       \qquad u_q\ne0.
 \label{eq:charge-factorization}
\end{equation}
Here $u_q$ is built from the nonzero Gaussian coefficient, Riesz
conjugation, and the nonzero functional-equation orientation.  Consequently
\begin{equation}
                  g_G=0\quad\Longleftrightarrow\quad\text{RH}.
 \label{eq:green-iff-rh}
\end{equation}
Equation \eqref{eq:green-iff-rh} is a detector theorem, not a proof of RH.
Its value is that it identifies the exact unsquared atomic coefficient to
be killed.

Resolved provenance retains the pair
\begin{equation}
                              \Delta_G=(g_G,-g_G).
 \label{eq:antidiagonal}
\end{equation}
The codiagonal is zero unconditionally:
\[
                              g_G+(-g_G)=0.
\]
The resolved pair itself vanishes exactly under RH.  Existing off-axis
cofaces saturate precisely the zero-provenance hyperplane, while the quotient
retains the incidence provenance.  Hence present cofaces cannot annihilate a
nonzero $\Delta_G$; direct augmentation merely forgets it.

This is the structural reason the proof must remain before trace,
positivity, and reflection-invariant quotienting.  Those consumers see the
zero codiagonal or a square of $D_\rho$.  The required conclusion concerns
the signed summand.

\section{The exact remaining comparison}

The two established halves are now:
\begin{equation}
 \begin{array}{ccl}
 \text{arithmetic source} &:&
       \omega_p\rightsquigarrow\tau_p\ne0,\\[1mm]
 \text{Riemann transport} &:&
       C_H\rightsquigarrow u_G\rightsquigarrow
       g_G(q)=D_{\rho_q}u_q.
 \end{array}
 \label{eq:two-halves}
\end{equation}
The missing mathematics is neither another source nor another equivalence
with RH.  It is a chain-level comparison between these lines.

\begin{definition}[Arithmetic comparison theorem]
\label{def:comparison}
An arithmetic comparison consists of source-defined maps, formed before
selecting a Xi zero, whose evaluation at each divisor copy $q$ and prime
$p$ gives a relation-kernel character $\mathfrak T_{p,q}(g_G)$ and has the
following properties.
\begin{enumerate}[label=(C\arabic*)]
\item \textbf{Coefficient fidelity:}
\begin{equation}
                    \mathfrak T_{p,q}(g_G)=D_{\rho_q}\tau_p.
 \label{eq:comparison-factor}
\end{equation}
\item \textbf{Chain naturality:} the construction commutes with the
Calder\'on--Witt cup differential, the free rational cover, harmonic Green
transport, and the prime--Gamma/co-Poisson boundary map.
\item \textbf{Literal boundary exactness:} the source representative of
$g_G$ is carried to the already annihilated Meyer--Poisson boundary sector,
so
\begin{equation}
                           \mathfrak T_{p,q}(g_G)=0.
 \label{eq:comparison-zero}
\end{equation}
\item \textbf{Source independence:} the maps commute with finite
rational-Witt refinement and prime exhaustion; no zero-dependent coefficient
is inserted into their definition.
\end{enumerate}
\end{definition}

Conditions (C1)--(C4) are not four unrelated conjectures.  They are the
coefficient, differential, boundary, and naturality faces of one chain map.
The source complex already has a secondary regulator whose exact kernel is
the reflected prime graph and which annihilates every literal
Meyer--Poisson boundary.  What is not known is that the completed Green
provenance comparison lands in that literal boundary complex with
coefficient \eqref{eq:comparison-factor}.

\begin{theorem}[Conditional capstone]
\label{thm:conditional-capstone}
If an arithmetic comparison in the sense of
Definition~\ref{def:comparison} exists, then the Riemann hypothesis holds.
\end{theorem}

\begin{proof}
Fix a divisor copy $q$ and a prime $p$.  By
\eqref{eq:comparison-factor} and \eqref{eq:comparison-zero},
\[
                              D_{\rho_q}\tau_p=0.
\]
Theorem~\ref{thm:calibrator} gives $k\in K_p$ with $\tau_p(k)\ne0$.
Evaluating at $k$ and using that $S_p$ is a complex vector space gives
$D_{\rho_q}=0$.  This holds for every multiplicity copy $q$, so every
nontrivial zero has real part $1/2$.
\end{proof}

The proof is deliberately short.  All difficulty is localized in the
comparison identity, not hidden in the final implication.

\subsection{Why this bridge is not circular}

A valid construction must pass four anti-circularity tests.
\begin{enumerate}
\item $\tau_p$ is defined and proved nonzero before any divisor is selected.
\item Its nonvanishing comes from the zero radical of the vector-valued
source cup, not from assuming an off-line zero.
\item Boundary vanishing occurs in the literal Meyer--Poisson complex, not
by imposing $g_G=0$ in the completed divisor carrier.
\item Coefficient fidelity retains $D_\rho$, not $D_\rho^2$ or a sum over a
functional-equation orbit.
\end{enumerate}
An RH-equivalent norm bound, a self-adjoint operator whose construction
already assumes criticality, or a divisor-indexed interpolation would fail
these tests.

\section{What is closed and what remains open}

\begin{longtable}{>{\raggedright\arraybackslash}p{0.29\textwidth}
                  >{\raggedright\arraybackslash}p{0.18\textwidth}
                  >{\raggedright\arraybackslash}p{0.43\textwidth}}
\toprule
Object or assertion & Status & Exact content \\
\midrule
\endfirsthead
\toprule
Object or assertion & Status & Exact content \\
\midrule
\endhead
Literal completed source & Closed & Fixed two-line prime--Gamma Cauchy
boundary is Bochner integrable and lands on one nonzero reflected Gaussian
horizon.\\
Harmonic primitive & Closed & $W=i\operatorname{sign}(D)$ satisfies
$W^2=-I$ and gives the unique completed Green primitive.\\
Completion wall & Closed & Compact Green range is exactly finite jets;
horizon lies in its closure but not its range; completed Green reflection is
the unique isometric extension.\\
Independent adaptive source & Closed & Adaptive M\"obius--theta columns
converge to the same primitive and Green detector.\\
Resolved provenance & Closed & Actual defect is antidiagonal; codiagonal is
zero unconditionally; resolved vanishing is equivalent to RH.\\
Arithmetic cup & Closed & Calder\'on--Witt cup is alternating, refinement
stable, and has zero vector-valued radical.\\
Relation-kernel calibrator & Closed & Free rational exactification produces
a canonical nonzero $\tau_p$ for every prime.\\
Literal boundary regulator & Closed on its stated domain & Secondary
Meyer--Poisson regulator annihilates every literal boundary and has exact
reflected-graph radical.\\
Green-to-kernel comparison & Open & Construct one source-defined chain map
giving $\mathfrak T_{p,q}(g_G)=D_{\rho_q}\tau_p$ and landing in the literal
boundary kernel.\\
Riemann hypothesis & Open & Follows immediately from that comparison, but
is not asserted here.\\
\bottomrule
\end{longtable}

This ledger rules out a common misdiagnosis: several source and transport
routes are open, but they do not have unrelated final obstacles.  Cauchy,
adaptive, Hodge, incidence, and resolved-provenance realizations already
meet at the same $g_G$.  Their shared missing object is the chain comparison
that turns literal arithmetic boundary exactness into a statement about the
retained relation-kernel character.

\section{No-go theorems explained by the same frontier}

The relation-kernel viewpoint clarifies why many correct constructions did
not produce zero exclusion.

\subsection{Codiagonal erasure}

The pair $(g_G,-g_G)$ has zero total for every divisor, on or off the line.
Any consumer that first applies the codiagonal cannot distinguish RH from a
nonzero antidiagonal.  This includes direct trace-like scalarization of the
two factor defects.

\subsection{Positive squaring}

Positive forms naturally see $|D_\rho|^2$ or a positive weighted sum.  They
can detect noncriticality after it occurs, but they do not supply a signed
source identity forcing each $D_\rho$ to vanish.  The calibrator $\tau_p$ is
linear and unsquared.

\subsection{Compact completion}

Finite Green jets are dense but miss the Gaussian horizon.  Ordinary
compact-range completion cannot create the missing vector; the forced
Hilbert completion does.  Once completed, however, it still only transports
$D_\rho$.  Functional analysis solves landing, not arithmetic annihilation.

\subsection{Balanced finite cofaces}

Finite balanced cofaces lie transverse to the nonzero antidiagonal, and
the off-axis closed span of existing cofaces is precisely the zero-provenance
hyperplane.  Adding more cofaces of the same type cannot kill the retained
provenance class.

\subsection{Source-only Koszul resolution}

All-prime label Koszul complexes resolve the homogeneous source directions
but leave a scalar Xi coefficient.  Higher homogeneous source resolution
therefore cannot remove horizontal charge unless prime-cycle data are
entangled with divisor provenance.  Equation \eqref{eq:comparison-factor}
is exactly such an entanglement, with a source-defined nonzero calibrator.

\subsection{Fixed carriers and bounded graded operators}

A fixed nonzero carrier multiplied by $D_\rho$ detects the charge but cannot
annihilate it independently.  Likewise, reflected valuation sectors may be
finite-order on one side and flat beyond every finite Rees order on the
other; finite-Rees-moderate or bounded graded sewing preserves that flat
radical.  The relation-kernel comparison is allowed to be singular at the
source boundary, but its arithmetic definition and coefficient fidelity
must remain exact.

\section{Why the \texorpdfstring{$8\pi$}{8-pi} laws do not close the proof}

The unweighted Ward energy obeys a sharp phase law: it is zero under RH and,
if RH fails, lies strictly above $8\pi$.  A different, endpoint-weighted Ward
index equals $2\pi$ times the total off-line endpoint charge, is strictly
below $0.093\pi$, and also vanishes exactly under RH.  These statements do
not contradict one another and cannot be combined into RH.

The measures differ.  Every unweighted off-line packet carries the same
$8\pi$ quantum, whereas its endpoint weight depends on height and has no
uniform positive lower bound.  A hypothetical high zero can therefore have
full unweighted Ward cost and arbitrarily small weighted cost.  The missing
comparison cannot be replaced by comparing the two totals.  It must act
atomically before weighting, which is precisely the role of
$D_\rho\tau_p$.

\section{Separation from classical routes}

The following comparison is structural.  It makes no priority claim and no
claim that every conceivable realization of a named program is impossible.

\begin{longtable}{>{\raggedright\arraybackslash}p{0.18\textwidth}
                  >{\raggedright\arraybackslash}p{0.30\textwidth}
                  >{\raggedright\arraybackslash}p{0.42\textwidth}}
\toprule
Program & Classical proof object & Present route \\
\midrule
\endfirsthead
\toprule
Program & Classical proof object & Present route \\
\midrule
\endhead
Hilbert--P\'olya & Construct a self-adjoint operator whose spectral data
encode zeros. & No spectral realization is required.  Horizontal purity is
to be forced before any operator spectrum exists, by vanishing of a
relation-kernel transgression.\\
Weil positivity & Prove positivity of an explicit-formula quadratic form. &
Retain the unsquared atomic coefficient $D_\rho$ and compare it with a
nonzero vector-valued arithmetic character.  Positivity is downstream.\\
Trace and cyclic programs & Realize zeros through a scaling flow, trace
formula, or cyclic quotient. & Work before trace and codiagonalization,
where $(g_G,-g_G)$ is still resolved.  Trace consumers may receive the
result but cannot supply the first vanishing arrow.\\
Sonin/Calkin and quotient-spectral programs & Detect zeros through
subspaces, quotient operators, essential spectrum, or boundary classes. &
Treat those classes as downstream images.  The decisive invariant lives on
the relation kernel of the source cover and must be compared before quotient
radicals enlarge.\\
Deninger-style cohomology & Seek Frobenius flow, Hodge theory, and purity
whose spectral determinant is zeta. & Use the Hodge and Green language, but
replace a hypothetical spectral cohomology by a constructed arithmetic cup,
free rational cover, and one explicit transgression comparison.\\
Explicit-formula scalarization & Sum prime and zero data against test
functions. & Keep multiplicity copies and factor provenance until the
atomic identity $D_\rho\tau_p=0$ is obtained.\\
\bottomrule
\end{longtable}

Connes's trace realization and Meyer's adelic/spectral construction are
representative downstream spectral programs in the sense used in this
table \cite{Connes1999,Meyer2005}.

The route is therefore pre-spectral, cohomological, and source-pointed.  It
does not refute Hilbert--P\'olya in full generality.  It shows why the current
source chain does not need Hilbert--P\'olya and why a spectral construction
would be downstream of the missing arithmetic theorem rather than a
substitute for it.

\section{Research program forced by the theorem}

The next work should not produce another RH equivalent.  It should construct
the comparison in Definition~\ref{def:comparison}.  The shortest credible
program has four concrete stages.

\begin{enumerate}
\item \textbf{Common chain object.}  Place the oriented prime--Gamma Cauchy
boundary and the Calder\'on--Witt first suspension in one source complex.
The chain object must retain prime label, valuation depth, rational-Witt
packet, reflection orientation, and Green provenance.
\item \textbf{Transgression naturality.}  Prove that passing to the free
rational cover commutes with the singular Gamma/co-Poisson boundary.  This
must identify the primitive restriction on $K_p$, not merely its scalar
trace.
\item \textbf{Coefficient computation.}  Evaluate the comparison on the
fixed Gaussian source and prove that its only divisor-dependent scalar is
the already established unsquared fracture $D_\rho$.  The remaining factor
must be exactly $\tau_p$, up to one fixed nonzero normalization independent
of $\rho$.
\item \textbf{Literal annihilation.}  Show that the compared Green class is
represented by an actual Meyer--Poisson boundary in the domain of the
secondary regulator.  Existing boundary annihilation then supplies
\eqref{eq:comparison-zero}.
\end{enumerate}

Failure at any stage has a useful positive survivor.  Failure of stage one
identifies a missing incidence coordinate; failure of stage two identifies
a genuine anomaly in Gamma reflection; failure of stage three produces a
new arithmetic coefficient besides horizontal charge; failure of stage four
identifies the precise quotient between the Cauchy and Meyer boundary
complexes.  None should be hidden by adding another scalar detector.

\section{Conclusion}

An exactifying cover does not destroy a degree-two class.  Under degree-one
descent it moves that class canonically to a character on the relation
kernel, and an alternating rational class vanishes precisely when that
character vanishes.  The Calder\'on--Witt source realizes this law with a
nonzero divisor-independent calibrator $\tau_p$.

The Riemann source and transport problem has meanwhile narrowed to one
oriented coefficient.  A fixed literal prime--Gamma boundary lands on a
nonzero Gaussian horizon; physical harmonic Hodge gives its unique completed
Green primitive; adaptive M\"obius--theta columns reach the same object; and
resolved incidence retains the Green defect
$g_G(\rho)=D_\rho u_\rho$ with $u_\rho\ne0$.  Direct codiagonalization kills
the pair $(g_G,-g_G)$ without killing either component.  Positivity, bounded
completion, source-only resolution, and fixed-carrier transport consequently
detect or erase the defect but do not force it to vanish.

The remaining theorem is exact and atomic:
\[
 \boxed{
 \Trans(\text{completed Green provenance at }\rho)
       =D_\rho\tau_p,
 \qquad
 \Trans(\text{literal source boundary})=0.}
\]
Because $\tau_p\ne0$, this identity would force $D_\rho=0$ one zero at a
time and prove RH.  Establishing that chain comparison---not finding another
spectrum, positive form, source, completion, or RH equivalence---is the
frontier of this route.

\begin{thebibliography}{9}
\small

\bibitem{Connes1999}
A.~Connes,
\newblock Trace formula in noncommutative geometry and the zeros of the
Riemann zeta function,
\newblock \emph{Selecta Math. (N.S.)} \textbf{5} (1999), 29--106.
\newblock \url{https://doi.org/10.1007/s000290050042}.

\bibitem{Meyer2005}
R.~Meyer,
\newblock A spectral interpretation for the zeros of the Riemann zeta
function,
\newblock preprint, arXiv:math/0412277 (2004).
\newblock \url{https://arxiv.org/abs/math/0412277}.

\bibitem{Riemann1859}
B.~Riemann,
\newblock \emph{\"Uber die Anzahl der Primzahlen unter einer gegebenen
Gr\"osse},
\newblock Monatsber. K\"onigl. Preuss. Akad. Wiss. Berlin (1859), 671--680.

\bibitem{Titchmarsh1986}
E.~C. Titchmarsh,
\newblock \emph{The Theory of the Riemann Zeta-Function},
\newblock second ed., revised by D.~R. Heath-Brown, Oxford University Press,
1986.

\end{thebibliography}
\normalsize

\end{document}
